\documentclass[aspectratio=169]{beamer}

% --- Theme & Colors ---
\usetheme{Madrid}
\usefonttheme{serif}

\definecolor{baselBg}{HTML}{fafaf8}
\definecolor{baselText}{HTML}{1a1a2e}
\definecolor{baselAccent}{HTML}{e94560}

\setbeamercolor{structure}{fg=baselAccent}
\setbeamercolor{normal text}{fg=baselText, bg=baselBg}
\setbeamercolor{frametitle}{fg=baselText, bg=baselBg}
\setbeamercolor{title}{fg=baselText, bg=}
\setbeamercolor{subtitle}{fg=baselText!70}
\setbeamercolor{titlelike}{fg=baselText, bg=}
\setbeamercolor{block title}{fg=white, bg=baselAccent}
\setbeamercolor{block body}{fg=baselText, bg=baselAccent!8}
\setbeamercolor{item}{fg=baselAccent}
\setbeamercolor{footline}{fg=baselText!50, bg=baselBg}

% --- Navigation & Footline ---
\setbeamertemplate{navigation symbols}{}
\setbeamertemplate{footline}{%
  \hfill\insertframenumber\kern1em\vskip4pt%
}

% --- Red underline on frame titles, with top padding ---
\addtobeamertemplate{frametitle}{\vskip0.4em}{%
  \vskip-0.3em%
  {\color{baselAccent}\rule{\textwidth}{0.8pt}}%
  \vskip0.3em%
}

% --- Packages ---
\usepackage{slidesonnet}
\usepackage{amsmath,amssymb}

% --- Bernoulli quote environment ---
\usepackage{tcolorbox}
\newtcolorbox{bquote}{%
  colback=baselAccent!10,
  colframe=baselAccent,
  leftrule=4pt,
  rightrule=0pt,
  toprule=0pt,
  bottomrule=0pt,
  arc=0pt,
  boxsep=10pt,
  left=8pt,
  fontupper=\itshape\color{baselText},
  width=0.85\textwidth,
  center
}

% --- Title info ---
\title{The Basel Problem}
\subtitle{A story of an impossible sum, a young genius,\\and a surprise visit from $\pi$}
\author{}
\date{}

\begin{document}

% ============================================================
% Frame 1: Title (1 visual state, 1 sub-slide)
% ============================================================
\begin{frame}[plain]
  \titlepage
  \say{[warmly] היום אני רוצה לספר לכם את אחד הסיפורים האהובים עליי במתמטיקה: נדבר על בעיית בָּאזֶל --- [building suspense] בעיה שנותרה ללא פתרון במשך שמונים וחמש שנה, עד שמתמטיקאי צעיר פיצח אותה עם רעיון פשוט ויפה.}
\end{frame}

% ============================================================
% Frame 2: A Simple Question (1 visual state, 1 sub-slide)
% ============================================================
\begin{frame}{A Simple Question}
  What happens when you add up the reciprocals of the perfect squares?

  \[
    \frac{1}{1} + \frac{1}{4} + \frac{1}{9} + \frac{1}{16}
    + \frac{1}{25} + \cdots
  \]

  In modern notation:

  \[
    \sum_{n=1}^{\infty} \frac{1}{n^2} = \ ?
  \]

  \say{הנה השאלה. קחו את כל המספרים השלמים בריבוע --- אחד, ארבע, תשע, שש עשרה, עשרים וחמש, וכן הלאה --- הִפכו את המספרים, וחָבְּרוּ את כולם. ... האם הסכום האינסופי הזה מתכנס? ואם כן --- [leaning in] לאיזה ערך?}
\end{frame}

% ============================================================
% Frame 3: The Challenge (1650) (2 visual states, 3 sub-slides)
% ============================================================
\begin{frame}{The Challenge (1650)}
  \textbf{Pietro Mengoli}, a mathematician in Bologna, posed this problem in 1650.

  \medskip
  The sum converges (by comparison with $\int_1^\infty x^{-2}\,dx = 1$).
  The partial sums creep toward something near \textbf{1.6449}\ldots

  \medskip
  But \emph{what} is this number, exactly?

  \say{הבעיה הוצגה לראשונה על ידי Pietro מֶנְגוֹ֫לִי בשנת אלף שש מאות וחמישים. הסכום אכן מתכנס --- אפשר לבדוק את זה בהשוואה לאינטגרל של אחד חלקי x בריבוע. חָבְּרוּ את מאה האיברים הראשונים, ותקבלו בערך 1 נקודה 6 3 4 9. חָבְּרוּ אלף איברים? תקבלו מעט יותר. הסכום בבירור מתקרב למשהו ליד 1 נקודה 6-4-4-9. [frustrated] אבל מהו המספר הזה? שורש של פולינום כלשהו? לוֹגַרִ֫יתְם של משהו? אף אחד לא הצליח לפענח את זה.}

  \pause

  \bigskip
  Even Jakob Bernoulli, one of the greatest mathematicians of his era,
  publicly admitted defeat in 1689:

  \say[2]{[gravely] אפילו בֶּרְנוּ֫לִי, אחד המתמטיקאים הגדולים של תקופתו, הביע את תסכולו בשנת אלף שש מאות שמונים ותשע. הוא כתב:}

  \pause

  \begin{bquote}
    ``If anyone finds and communicates to us that which thusfar
    has eluded our efforts, great will be our gratitude.''

    \upshape\raggedleft--- Jakob Bernoulli, 1689
  \end{bquote}

  \say[3, voice=bernoulli]{"If anyone finds and communicates to us that which thusfar has eluded our efforts, great will be our gratitude."}
\end{frame}

% ============================================================
% Frame 4: Enter Leonhard Euler (2 visual states, 2 sub-slides)
% ============================================================
\begin{frame}{Enter Leonhard Euler}
  The problem remained open for \textbf{85 years}.

  \medskip
  Then, in 1734 in Saint Petersburg, a young mathematician from Basel
  named \textbf{Leonhard Euler} found the answer.

  \say{במשך שמונים וחמש שנה, אף אחד לא הצליח לפתור את הבעיה. ... ואז [brightening] בשנת אלף שבע מאות שלושים וארבע, מתימטיקאי בן עשרים ושבע בשם לֵאוֹנָארְד אוֹיְלֶר, שעבד באקדמיה למדעים בסנט פטרסבורג, מצא את התשובה. הוא הציג את ההוכחה שלו בשנה שלאחר מכן.}

  \pause

  \bigskip
  \[
    \sum_{n=1}^{\infty} \frac{1}{n^2} = \frac{\pi^2}{6}
  \]

  \medskip
  $\pi$? In a sum that has nothing to do with circles?

  \say[2]{הסכום שווה ל-פאי בריבוע חלקי שש. [shocked] וזה היה מפתיע. פאי הוא היחס בין היקף המעגל לקוטרו --- למה הוא מופיע בסכום של הופכי ריבועים?! אין כאן מעגל בשום מקום! [calming down] בואו אראה לכם את הטיעון המקורי של אוֹיְלֶר. הוא לא לגמרי מדויק לפי סטנדרטים מודרניים, אבל האינטואיציה כל כך משכנעת שכולם קיבלו אותה מיד, ואת הפערים בהוכחה סגרו במהלך המאה שלאחר מכן.}
\end{frame}

% ============================================================
% Frame 5: Step 1 — sin(x) (3 visual states, 3 sub-slides)
% ============================================================
\begin{frame}{Step 1: Start with $\sin(x)$}
  We know the Taylor series for sine:

  \[
    \sin(x) = x - \frac{x^3}{3!} + \frac{x^5}{5!} - \frac{x^7}{7!} + \cdots
  \]

  \say{ההוכחה של אוֹיְלֶר מתחילה עם משהו שאולי לא ציפיתם לו --- [intrigued] הפונקציה סִינוּס. אנחנו יודעים את טור טֵיְלוֹר של סִינוּס של x: זה x פחות x בשלישית חלקי שלוש עצרת, ועוד x בחמישית חלקי חמש עצרת, וכן הלאה.}

  \pause

  \medskip
  Divide both sides by $x$:

  \[
    \frac{\sin(x)}{x} = 1 \alt<3>{{\color{baselAccent}- \frac{x^2}{6}}}{- \frac{x^2}{6}} + \frac{x^4}{120} - \cdots
  \]

  \say[2]{עכשיו זכרו, טור אינסופי הוא פשוט גבול של סכומים חלקיים. לכל x קבוע ששונה מאפס, אפשר לחלק את הגבול ב-x --- זה מוצדק על ידי משפט הגבול האלגברי. אז מקבלים סִינוּס של x חלקי x שווה אחד פחות x בריבוע חלקי שש, ועוד איברים מסדר גבוה יותר.}

  \pause

  \bigskip
  Note the coefficient of $x^2$ on the right: it is ${\color{baselAccent}-\frac{1}{6}}$.

  \say[3]{שימו לב למקדם של x בריבוע --- מינוס שישית. [knowingly] המספר הזה עוד יחזור!}
\end{frame}

% ============================================================
% Frame 6: Step 2 — Factor by roots (2 visual states, 2 sub-slides)
% ============================================================
\begin{frame}{Step 2: Factor by the Roots}
  Euler's bold move: $\dfrac{\sin(x)}{x}$ has zeros at
  $x = \pm\pi,\; \pm 2\pi,\; \pm 3\pi,\; \ldots$

  \say{[with excitement] ועכשיו הצעד הנועז של אוֹיְלֶר --- הצעד שהפך את ההוכחה הזו למפורסמת. הפונקציה סִינוּס של x חלקי x מתאפסת כאשר x שווה פלוס מינוס פאי, פלוס מינוס שני פאי, פלוס מינוס שלוש פאי, וכן הלאה.}

  \pause

  \bigskip
  Euler treated it like a polynomial and factored it by its roots:

  \[
    \frac{\sin(x)}{x}
    = \left(1 - \frac{x^2}{\pi^2}\right)\!
    \left(1 - \frac{x^2}{4\pi^2}\right)\!
    \left(1 - \frac{x^2}{9\pi^2}\right)\cdots
  \]

  \medskip
  \small
  \textbf{Assumes:} An entire function is determined by its zeros
  (up to normalization). Justified by the Weierstrass factorization
  theorem (1876).

  \say[2]{אוֹיְלֶר חשב בְּאנלוגיה. אם לפולינום יש שורשים מסוימים, הרי שאפשר לכתוב אותו כמכפלה של גורמים, אחד לכל שורש. הוא החיל את אותו היגיון על סִינוּס של x חלקי x --- [raising eyebrows] למרות שזה בכלל לא פולינום. הוא כתב את הביטוי כמכפלה אינסופית, כאשר כל גורם מתאים לזוג שורשים. [carefully] עכשיו, זו הייתה "הנחה יצירתית", לא הוכחה. ההצדקה הפורמלית הגיעה הרבה יותר מאוחר, עם משפט הפירוק של וַאיֶרְשְׁטְרָס מ-1876. אז מה זה נותן לנו?}
\end{frame}

% ============================================================
% Frame 7: Step 3a — x^2 from a product (1 visual state, 1 sub-slide)
% ============================================================
\begin{frame}{Step 3a: How to Read the $x^2$ Coefficient from a Product}
  Consider just \textbf{three} factors --- write $a_k = \frac{1}{k^2\pi^2}$
  for short:

  \begin{gather*}
    (1 - a_1 x^2)(1 - a_2 x^2)(1 - a_3 x^2) \\[4pt]
    = 1 - (a_1 + a_2 + a_3)\,x^2 + (\text{terms in } x^4, x^6)
  \end{gather*}

  \medskip
  Each $x^2$ contribution comes from picking $-a_k x^2$ from \textbf{one}
  factor and~$1$ from the rest. Cross-terms yield $x^4$ or higher.

  \medskip
  So the $x^2$ coefficient is just $-(a_1 + a_2 + a_3)$.

  \say{לפני שניגש למכפלה האינסופית, בואו נראה מה קורה עם מכפלה סופית. קחו רק שלושה גורמים. כל גורם נראה כמו אחד פחות קבוע כפול x בריבוע.
    כשמכפילים, המקדם של x בריבוע הוא פשוט סכום המקדמים בגורמי המכפלה. עיצרו רגע את הסרטון כדי לוודא שאתם רואים למה זה נכון.}
\end{frame}

% ============================================================
% Frame 8: Step 3b — Infinite product (1 visual state, 1 sub-slide)
% ============================================================
\begin{frame}{Step 3b: Apply to the Infinite Product}
  \[
    \prod_{n=1}^{\infty}\!\left(1 - \frac{x^2}{n^2\pi^2}\right)
    = 1 - \left(\sum_{n=1}^{\infty}\frac{1}{n^2\pi^2}\right)x^2 + \cdots
  \]

  \bigskip
  The $x^2$ coefficient is:

  \[
    -\left(\frac{1}{\pi^2} + \frac{1}{4\pi^2}
    + \frac{1}{9\pi^2} + \cdots\right)
    = -\frac{1}{\pi^2}\sum_{n=1}^{\infty}\frac{1}{n^2}
  \]

  \medskip
  \small
  \textbf{Assumes:} Absolute convergence and term-by-term expansion,
  as for finite products.

  \say{עכשיו נרחיב את זה למכפלה האינסופית. בדיוק לפי אותו היגיון, המקדם של x בריבוע הוא סכום כל ה-מקדמים במכפלה עם סימן מינוס. זה נותן לנו מינוס אחד חלקי פאי בריבוע, ועוד אחד חלקי ארבע פאי בריבוע, ועוד אחד חלקי תשע פאי בריבוע, וכן הלאה. הוציאו את אחד חלקי פאי בריבוע כגורם משותף, ומה נשאר? [triumphantly] בדיוק הסכום המסתורי שלנו. [more measured] עכשיו, כדי להרחיב את הטיעון ממכפלה סופית לאינסופית, צריך לוודא שזה בכלל מותר לנו. זה נכון כאן, אבל כלל לא הוכח בתקופתו של אוֹיְלֶר.}
\end{frame}

% ============================================================
% Frame 9: Step 4 — The Punchline (3 visual states, 3 sub-slides)
% ============================================================
\begin{frame}{Step 4: The Punchline}
  From the Taylor series:

  \[
    \frac{\sin(x)}{x}
    = 1 \mathbin{{\color{baselAccent}- \frac{1}{6}}}\, x^2 + \cdots
  \]

  \say{עכשיו נשווה. טור טֵיְלוֹר אמר לנו שהַמְּקַדֵּם של x בריבוע הוא מינוס שישית.}

  \pause

  \medskip
  From the product:

  \[
    \frac{\sin(x)}{x}
    = 1 \mathbin{{\color{baselAccent}
          - \frac{1}{\pi^2}\sum_{n=1}^{\infty}\frac{1}{n^2}}}\, x^2 + \cdots
  \]

  \say[2]{והמכפלה האינסופית אומרת לנו שזה מינוס אחד חלקי פאי בריבוע כפול הסכום שלנו.}

  \pause

  \bigskip
  Coefficients must agree:
  $\quad\dfrac{1}{\pi^2}\displaystyle\sum_{n=1}^{\infty}\frac{1}{n^2}
    = \frac{1}{6}$

  \say[3]{שני הביטויים הם טורי חזקות של אותה פונקציה, ולפונקציה יש רק ייצוג אחד כטור חזקות. אז המקדמים חייבים להיות שווים. הַשְׁווּ אותם, וקבלו: [deliberately] אחד חלקי פאי בריבוע כפול הסכום שווה שישית.}
\end{frame}

% ============================================================
% Frame 10: Result (1 visual state, 1 sub-slide)
% ============================================================
\begin{frame}{$\displaystyle\sum_{n=1}^{\infty}\frac{1}{n^2}
      = \frac{\pi^2}{6}$}
  Multiply both sides by $\pi^2$:

  \[
    1 + \frac{1}{4} + \frac{1}{9} + \frac{1}{16}
    + \frac{1}{25} + \cdots = \frac{\pi^2}{6} \approx 1.6449\ldots
  \]

  \say{הַכפילו את שני הצדדים ב-פאי בריבוע, והנה זה. אחד ועוד רבע ועוד תשיעית ועוד אחד חלקי שש עשרה, וכן הלאה עד אינסוף --- שווה פאי בריבוע חלקי שש. [in awe] אותו 1 נקודה 6 4 4 9 מסתורי שאף אחד לא הצליח לזהות במשך שמונים וחמש שנה? ... פאי בריבוע חלקי שש. [softly, with wonder] קבוע המעגל הסתתר בתוך הסכום כל הזמן...}
\end{frame}

% ============================================================
% Frame 11: Why It Matters (2 visual states, 2 sub-slides)
% ============================================================
\begin{frame}{Why It Matters}
  Euler's proof was \textbf{not rigorous} by modern standards --- the
  factoring step needed Weierstrass's theory of infinite products to
  justify, over 140 years later (1876).

  \medskip
  But the result itself was correct, and it opened the door to a vast
  landscape.

  \say{אז האם צריך לדאוג שההוכחה של אוֹיְלֶר דילגה על שלבים? [reassuringly] מתמטיקאים בהחלט דאגו. עברו יותר מ*מאה שנה* עד שוַאיֶרְשְׁטְרָס סיפק את הכלים שהפכו את שלב הפירוק למוצדק לחלוטין. אבל התשובה מעולם לא הייתה מוטלת בספק. [reflectively] לפעמים האינטואיציה הגיעה הרבה לפני ההוכחה --- ובפער הזה נולדה מתמטיקה חדשה.}

  \pause

  \bigskip
  The sum we computed is $\zeta(2)$, a special value of the
  \textbf{Riemann zeta function}:

  \[
    \zeta(s) = \sum_{n=1}^{\infty} \frac{1}{n^s}
  \]

  \begin{itemize}
    \item Euler showed all \textbf{even} values
          $\zeta(2), \zeta(4), \zeta(6), \ldots$ are rational multiples
          of powers of $\pi$

          \say[2]{מה שחישבנו הוא למעשה ערך מיוחד של מה שנקרא היום פונקציית זֶ֫טָא של רִימָן --- עבור-s שווה שתיים. אוֹיְלֶר *עצמו* המשיך והראה שכל הערכים הזוגיים, זֶ֫טָא של שתיים, זֶ֫טָא של ארבע, זֶ֫טָא של שש, וכן הלאה, הם כפולות רציונליות של חזקות של פאי.}

          \pause

    \item But the \textbf{odd} values
          $\zeta(3), \zeta(5), \zeta(7), \ldots$ remain mysterious ---
          no closed forms are known
  \end{itemize}

  \say[3]{אבל הערכים האי-זוגיים? זֶ֫טָא של שלוש, זֶ֫טָא של חמש, זֶ֫טָא של שבע? *עדיין* לא ידועים... עד היום, אף אחד לא מצא ביטויים סגורים לאף אחד מהם. כמה מהבעיות הפתוחות הגדולות ביותר במתמטיקה חוזרות בדיוק אל הסכום היפה שלנו.}
\end{frame}

% ============================================================
% Frame 12: What We Learned (1 visual state, 1 sub-slide)
% ============================================================
\begin{frame}{What We Learned}
  \[
    1 + \frac{1}{4} + \frac{1}{9} + \frac{1}{16} + \cdots
    = \frac{\pi^2}{6}
  \]

  \begin{itemize}
    \item A simple question can hide a profound answer
    \item The boldness to treat $\sin(x)/x$ as a polynomial was the
          key insight
    \item $\pi$ connects geometry and number theory in ways we are
          still discovering
  \end{itemize}

  \say{[warmly] אז מה למדנו? לפעמים שאלה פשוטה --- כזו שאפשר להסביר לילד קטן --- מסתירה תשובה עמוקה להפליא. הגאונות של אוֹיְלֶר לא הייתה רק טכנית. הייתה לו תעוזה לנסות משהו שלא היה אמור לעבוד וחוש חד לאתר יופי מתמטי. [with wonder] ובלב של הכל נמצא פאי, שמחבר את הגאומטריה של מעגלים, לאריתמטיקה של מספרים שלמים בְּצוּרוֹת שאנחנו עדיין מגלים היום. [concluding] תודה! שהצטרפתם אליי למסע הזה דרך אחת ההוכחות היפות ביותר במתמטיקה... נתראה בסרטון הבא!}
\end{frame}

% ============================================================
% Frame 13: Sources (1 visual state, silent)
% ============================================================
\begin{frame}{Sources}
  \small
  \begin{itemize}
    \item P.\ Mengoli, \emph{Novae quadraturae arithmeticae, seu de
            additione fractionum} (Bologna, 1650)
    \item J.\ Bernoulli, \emph{Tractatus de seriebus infinitis}
          (Basel, 1689) --- quote per W.\ Dunham,
          \emph{Euler: The Master of Us All} (1999)
    \item L.\ Euler, ``De summis serierum reciprocarum'' (E41),
          presented Dec.\ 5, 1735,
          \emph{Commentarii academiae scientiarum Petropolitanae} 7 (1740)
    \item K.\ Weierstrass, ``Zur Theorie der eindeutigen analytischen
          Funktionen,'' \emph{Abhandlungen der K\"oniglichen Akademie der
            Wissenschaften zu Berlin} (1876)
    \item R.\ Ap\'ery, ``Irrationalit\'e de $\zeta(2)$ et $\zeta(3)$,''
          \emph{Ast\'erisque} 61 (1979)
  \end{itemize}
  \nonarration[3]
\end{frame}

\end{document}
